% Capstone 3: index calculus. Port of «Мини-курсовая - метод Исчисления индексов».
% The performance numbers come from the generated requirements.tex (make bench).
\documentclass{ca-capstone}
\capstonerefs{alpertron-dilog,gmpy2,sagemath,hps}
\begin{document}
% Generated by `make bench CAPSTONE=index-calculus` (taskgen.capstone report) from
% bench/results.json: reference time x slack on the reference machine.
% Do not edit: change bench/config.yaml and re-run the benchmark.
\newcommand{\ReqSlack}{10}
\newcommand{\ReqBitsMin}{50}
\newcommand{\ReqBitsMax}{60}
\newcommand{\ReqSeconds}{50}

\capstonetitle{Решение задачи DLP методом Исчисления индексов}

\goal
Требуется реализовать метод Исчисления индексов для решения задачи дискретного
логарифмирования в $\GF{p} \cong \Zn{p}$ (т.\,е. метод решения сравнений вида
$g^x \equiv h \pmod{p}$).

\functionality
\begin{enumerate}
  \item Вычисление (1)~значения параметра~$B$, (2)~требуемого количества \emph{$B$-гладких}
    чисел и (3)~ожидаемого количества операций поиска сравнений при формировании системы на
    первом шаге метода\footnote{параметр~$B$ и требуемое количество итераций поиска
    выражаются в терминах функции~$L(p)$ (\emph{как?})}
  \item Поиск \emph{$B$-гладких} чисел и составление системы сравнений
  \item Решение системы сравнений для поиска значений дискретного логарифма по основанию~$g$
    простых чисел из фактор-базы
  \item Поиск \emph{$B$-гладкого} числа вида $h \cdot g^{-k} \pmod{p}$ и решение задачи
    \emph{DLP} с помощью найденных ранее дискретных логарифмов
\end{enumerate}

\requirements
Требования к реализации:
\begin{enumerate}
  \item корректность решения задачи \emph{DLP}
  \item приемлемая скорость работы (см. комментарий ниже)
\end{enumerate}

\textbf{Корректность работы} можно проверять, самостоятельно генерируя задачи, либо же
прибегая к помощи сторонних решателей, например, \cite{alpertron-dilog,sagemath}.

В данной работе важна не скорость, а корректность реализации и соответствие приведённым выше
требованиям. Однако тестовые примеры с числами размеров примерно
\textbf{\ReqBitsMin--\ReqBitsMax~бит} должны выполняться за приемлемое время "--- несколько
десятков секунд (не более \ReqSeconds~с).

\details
Так как цель работы "--- решение задачи \emph{DLP}, в своей реализации \textbf{можно
пользоваться библиотеками} для работы с большими целыми числами, а также пакетами линейной
алгебры.

\textbf{Проверку \emph{$B$-гладкости}} можно осуществлять напрямую: \textbf{путём пробных
делений} на все простые числа из фактор-базы.

Для \textbf{решения системы можно} воспользоваться функциями некоторого \textbf{пакета
линейной алгебры}.

Рекомендуется использовать язык \emph{Python} и систему компьютерной алгебры
\emph{SageMath}~\cite{sagemath}.

Желательно сохранять промежуточные результаты, такие как найденные на первом этапе
\emph{$B$-гладкие} степени, вид составленной системы сравнений и её решения, найденное на
втором этапе выражение и проч.

Платформо-независимый, \textbf{переносимый код} писать \textbf{не требуется}, но во время сдачи
должна быть возможность проверить работу.

\schedule
Работа выполняется \textbf{индивидуально}. Срок выполнения "--- \textbf{месяц}.

Работа должна быть ваша: реализация должна быть \textbf{разработана и написана вами}, а не
собрана из кусков чужих работ.

Желательно хранить код в \emph{git}-репозитории на \emph{github}, \emph{bitbucket} или подобных
платформах. \textbf{Документация не требуется}.

\remarks
\begin{enumerate}
  \item Оценка значений параметров приведена в разделе~3.7.1~\cite{hps}
\end{enumerate}

\references
\end{document}
